\section{Introduction}
\label{sec:intro}

Deploying deep neural networks directly on edge devices (e.g., smartphones, wearables, and microcontrollers) is essential to reduce inference latency, protect user privacy, and enable robust operation in disconnected environments. However, real-world edge devices are tightly constrained by physical limitations: restricted random-access memory (RAM), limited non-volatile storage, and strict thermal envelopes.

As edge applications grow in complexity, systems are increasingly required to perform multi-task learning (MTL), serving diverse functionalities simultaneously (e.g., scene classification, object detection, and segmenting multiple domains). Jointly training a single massive multi-task model from scratch is often practically infeasible due to data privacy constraints, immense computational costs, and susceptibility to catastrophic forgetting or task interference. 

To address these deployment challenges, \emph{model merging} has emerged as an exceptionally cost-effective paradigm \cite{wortsman2022model,ilharco2023editing}. Model merging combines independently fine-tuned task-specific experts into a single joint multi-task model without requiring original, massive training datasets or expensive retraining. By taking the difference between fine-tuned expert weights and their shared pre-trained base model, we can construct task vectors \cite{ilharco2023editing} and combine them using layer-wise coefficients \cite{yang2024adamerging}.

\begin{figure}[t]
  \centering
  \includegraphics[width=\columnwidth]{comparison_plot.png}
  \caption{Joint Mean Accuracy (\%) across four highly disparate visual tasks (MNIST, FashionMNIST, CIFAR-10, SVHN) using a ViT-Tiny backbone under various pruning pipelines. Our experiments expose a critical representational collapse: every merged model performs near random guessing (~10\% to 14\%), indicating a fundamental limitation in linear weight-space merging under high task conflict and low model capacity.}
  \label{fig:comparison_curves}
\end{figure}

Despite the elegance of model merging, a major pragmatic limitation remains: \textbf{the merged model is still fully dense (100\% parameter size).} To make model merging truly viable in the wild, we must compress the merged model to fit physical edge storage limits. Traditional weight pruning \cite{han2015learning} seems a natural candidate, but naively sequencing merging and pruning leads to severe suboptimality. 

To systematically investigate this intersection, we conduct a rigorous post-mortem and limitation-mapping study of joint model merging and pruning on resource-constrained edge hardware. We present and analyze \textbf{ZipMerge} (Post-Merge Joint Weight Pruning and Coefficient Tuning), a framework that integrates magnitude pruning directly into the test-time adaptation loop. Instead of treating merging and pruning as separate, decoupled stages, ZipMerge co-optimizes the continuous merging coefficients $\Lambda$ and the binary pruning mask $M$ simultaneously on a tiny, unlabeled calibration dataset (consisting of only 16 images per task) using an unsupervised minimum entropy loss. 

We explore two distinct optimization paradigms to solve this joint optimization problem:
\begin{enumerate}[leftmargin=*]
    \item \textbf{ZipMerge (STE):} Employs a Straight-Through Estimator (STE) \cite{bengio2013estimating} to propagate first-order gradients back to the continuous merging coefficients $\Lambda$ through the non-differentiable pruning mask operator.
    \item \textbf{ZipMerge (ES):} Employs a black-box 1+1 Evolution Strategy (1+1 ES) \cite{salimans2017evolution} to explore the coefficient space, sidestepping the non-differentiable step function of the pruning threshold.
\end{enumerate}

However, rather than presenting a curated narrative of triumph, we conduct a highly rigorous, honest empirical investigation using a compact \texttt{timm} Vision Transformer (\texttt{vit\_tiny\_patch16\_224}) backbone evaluated across four highly disparate datasets (MNIST, FashionMNIST, CIFAR-10, and SVHN). Our results expose severe empirical boundaries that represent a vital lesson for real-world practitioners:
\begin{itemize}[leftmargin=*]
    \item \textbf{Catastrophic Representational Collapse:} Under extreme domain and task shift, every single merged model, including Uniform, AdaMerging, and the proposed ZipMerge, completely collapses to random guessing levels (~10\% to 14\% absolute accuracy).
    \item \textbf{The Overfitting-Optimizer Paradox:} We reveal that unconstrained minimum-entropy TTA on tiny calibration sets overfits transductively, successfully minimizing entropy while destroying generalizable features and driving test-set accuracy down.
    \item \textbf{Prune-then-Merge Baseline Outperformance:} We show that the unoptimized, decoupled baseline, Prune-then-Merge (P-then-M), consistently outperforms test-time optimization because pre-merging pruning acts as an extreme regularizer, zeroing out conflicting parameter noise.
\end{itemize}

By studying and reporting these boundaries, we transition model merging from an academic toy scenario to the physical system realities, providing actionable architectural guidelines for edge engineers.
